\documentclass[aspectratio=169,usepdftitle=true,11pt,ngerman,t]{beamer}

\newcommand\lqfirstname{Name}
\newcommand\lqfullname{Name Last}
\newcommand\lqemail{name.last}
\errorcontextlines999999

% base
\usepackage[utf8]{inputenc}
\usepackage[T1]{fontenc}
\usepackage{lmodern}
\usepackage{microtype}
\usepackage[main=ngerman]{babel}
\usepackage{csquotes}

% graphics
\usepackage[prefix=]{xcolor-material}
\usepackage{graphics}
\usepackage{tcolorbox}
\usepackage{twemojis}
\usepackage{tikz}
\usetikzlibrary{
    positioning,
    calc,
    tikzmark,
    circuits.logic.IEC
}

% math
\usepackage{amsmath,amssymb}
\usepackage{mathtools}
\usepackage{array}
\usepackage{booktabs}
\usepackage{cancel}
\usepackage{ulem}

% flo
\usepackage{fancyqr}
\usepackage[addons]{color-palettes}
\usepackage[encoding,noloadlangs,cpalette,numinpar,fakeminted]{sopra-listings}
\usepackage{code-animation}
\usetheme[theorems]{colorful-dream}
\theoremstyle{exercise}
\newtheorem*{exercise*}{Aufgabe}
\theoremstyle{solve}
\newtheorem*{solve*}{Lösung}

% tools
\usepackage{algpseudocode}
\usepackage{dirtree}

% warnings
\usepackage{silence}
\WarningFilter{latexfont}{Some font shapes were not available}
\WarningFilter{latexfont}{Font shape}
\WarningFilter{pgf}{Returning node center instead of a point on node border}
\let\origenddocument\enddocument\renewcommand{\enddocument}{\hfuzz=10.95pt\origenddocument}
% title image: box
% ([vshift], inverse scale, width, content)
\newcommand\titleimagebox[4][0pt]{%
\newlength\tibheight\setlength\tibheight{4cm}
\newlength\tibtY\setlength\tibtY{-6em}
\titleimage{\begin{minipage}[c][#2\tibheight]{#3}\vspace*{#2\tibtY}%
    \vspace*{#1}\begin{tcolorbox}[boxrule=0.1pt]%
    \centering#4% TODO: bug: not centered if too wide
    \end{tcolorbox}%
\end{minipage}}}

% outro image: text
% (text)
\newcommand\outroimagetext[1]{\outro{Karlsruhe, \today%
\begin{tikzpicture}[remember picture,overlay]
    \node[align=center] at (current page.center) {#1};
\end{tikzpicture}}}% chktex 31

% reset counters
\newcommand\resetframecounters{\addtocounter{exercise}{-1}\addtocounter{solve}{-1}}
\newcommand\resetex{\addtocounter{exercise}{-1}}
\newcommand\resetsolve{\addtocounter{solve}{-1}}

% small footnote
\newcommand{\nicefootnote}[1]{\tikz[overlay,remember picture]{\node[anchor=south,text width=\dimexpr\textwidth+2em\relax,align=center,above=1em of current page.south] (nicefootnote) {\hrule\vspace{1pt}\footnotesize#1};}}

% kv
% ([draw], variables, values)
% {currently only works correctly for 4 variables}
\newcommand{\lqkv}[3][]{
    \begin{tikzpicture}
        \foreach \v [count=\k from 1] in {#2} {\global\let\Vdim=\k}
        \pgfmathsetmacro\rows{2^ceil(\Vdim/2)}
        \pgfmathsetmacro\cols{2^int(\Vdim/2)}

        \foreach \v [count=\k from 1] in {#2} {
            \pgfmathsetmacro\shift{0.5+floor(\k*ceil(\Vdim/2)/2)}

            \ifcase\k\relax
            \or
                \draw (\shift-1,0.75) -- (\shift+1,0.75);
                \node[fill=white] at (\shift,0.75) {\v};
            \or
                \draw (-0.75,-\shift) -- (-0.75,-\shift+2);
                \node[fill=white] at (-0.75,-\shift+1) {\v};
            \or
                \draw (\shift-2,-\cols+0.25) -- (\shift,-\cols+0.25);
                \node[fill=white] at (\shift-1,-\cols+0.25) {\v};
            \or
                \draw (\cols-0.25,-\shift+1) -- (\cols-0.25,-\shift+3);
                \node[fill=white] at (\cols-0.25,-\shift+2) {\v};
            \fi
        }
      
        \foreach \v [count=\k from 0] in {#3} {
            \pgfmathsetmacro\x{mod(\k,\rows)}
            \pgfmathsetmacro\y{int(\k/\rows)}
            \node[draw, minimum width=1cm, minimum height=1cm, align=center] at (\x,-\y) {\v};
        }

        \begin{scope}
            \clip (-0.5,0.5) rectangle (3.5,-3.5);
            \foreach \r in {#1} {\draw[rounded corners=1ex] \r;}
        \end{scope}

        \global\let\Vdim\undefined
    \end{tikzpicture}
}
\title{Digitaltechnik}
\institute{Karlsruher Institut für Technologie}

\author{\lqfullname}
\email{\lqemail @student.kit.edu}% chktex 1
% credit flo
\addtobeamertemplate{title page}{}{%
\scriptsize\color{gray}\begin{tikzpicture}[overlay, remember picture]
    \node[above left = 1em and 1em of current page.south east,align=center] (n0) {\LaTeX-Vorlage von\\\href{https://github.com/EagleoutIce}{Florian Sihler}};
    \node[above = 0cm of n0] (n1) {\fancyqr[size=1.5cm,tight]{https://eagleoutice.github.io/beamer-themes/\#colorful-dream}};
\end{tikzpicture}}% chktex 31

%%%%%%%%%%%%%%%%%%%%%%%%%%%%%%%%%%%%%%%%%%%%%%%%%%%%%%%%%%%%%%%%

\subtitle{Tutorium 2}
\date{19. Mai 2026}

\titleimagebox{1}{8cm}{\(1\) oder \(\overline{1}\), das ist hier die Frage!}

\outroimagetext{\begin{tabular}{c}
    \large Soon\textsuperscript{\texttrademark} auf \href{https://github.com/Laqco/kit-dt-ss26}{GitHub}!\\
    \tiny Wenn ich Zeit habe\dots
\end{tabular}}

%%%%%%%%%%%%%%%%%%%%%%%%%%%%%%%%%%%%%%%%%%%%%%%%%%%%%%%%%%%%%%%%

\title{Digitaltechnik}

\begin{document}

\section[\(k_1, k_2, m_1, \dots \)]{Hamming-Codes}
\subsection{Vektor}

\begin{frame}{Ein geheimer Code!}\onslide<+->
    \begin{exercise}
        \begin{enumerate}
            \item[1.] Wie viele Prüfbits werden für ein Datenwort der Länge 120 Bits benötigt?
        \end{enumerate}
    \end{exercise}\onslide<+->
    \begin{solve}
        \only<.>{\(2^k\geq m+k+1\) und \(2^{k-1}<m+(k-1)+1\)}
        \onslide<+->\(2^k\geq 121+k\) und \(2^{k-1}<120+k\)\onslide<+->\\
        \(\dots k=7\)
    \end{solve}
\end{frame}
\resetframecounters%

\begin{frame}{Ein geheimer Code!}\onslide<+->
    \begin{exercise}
        \begin{enumerate}
            \item[2.] Dekodieren Sie das folgende Hamming-Codewort: \T{101001100001101}
        \end{enumerate}
    \end{exercise}\onslide<+->
    \begin{solve}
        Prüfgleichungen: \(k_1: 1\quad k_2: 1\quad k_3: 1\quad k_4: 0\) \onslide<+->\\
        \(0111_2 = 7_{10}\) \onslide<+->\\
        \faAngleRight~\T{10100110\textcolor{red}{1}001101}
    \end{solve}
\end{frame}
\resetframecounters%

\begin{frame}{Ein geheimer Code!}\onslide<+->
    \begin{exercise}
        \begin{enumerate}
            \item[3.] Kodieren Sie die Zeichenkette \enquote{dt} mit dem Hamming-Code.
        \end{enumerate}
    \end{exercise}\onslide<+->
    \begin{solve}
        \enquote{d}: \onslide<+->\T{01100100}\\
        \enquote{t}: \onslide<+->\T{01110100}\onslide<+->\\
        \T{011000100011100101010}
    \end{solve}
\end{frame}

\subsection{Matrix}

\begin{frame}{The Matrix is everywhere.}\onslide<+->
    \begin{solve}
        \begin{enumerate}[\alph*)]% chktex 9 chktex 10
            \item<+-> \((1,0,1,1,0,1,0)\)
            \item<+-> \((0,1,1)\) \faAngleRight~Fehler an der 3. Stelle von \textbf{links}!\onslide<+->\\
                \(CW = (0,1,0,0,1,0,1)\)
            \item<+-> \dots
        \end{enumerate}
    \end{solve}
\end{frame}
\addtocounter{exercise}{1}

\section{Boole'sche Algebra}

\begin{frame}{Wieso gibt diese Aufgabe so viele Punkte?}\onslide<+->
    \alt<1>{\begin{exercise}
        \begin{enumerate}
            \item[1.] Um welche der folgenden Aussagen handelt es sich um eine Tautologie?
        \end{enumerate}
    \end{exercise}}{\begin{exercise}
        \begin{enumerate}
            \item[1.] Bei welchen der folgenden Aussagen handelt es sich um eine Tautologie?
        \end{enumerate}
    \end{exercise}}\pause%
    \begin{solve}
        \begin{itemize}[<+->]
            \item \((a\wedge b)\rightarrow a\)
            \item \((a\rightarrow b) \vee (b\rightarrow a)\)
            \item \((a\leftrightarrow b) \vee c\)
            \item \((a\wedge(b\vee c))\leftrightarrow ((a\wedge b)\vee(a\wedge c))\)
        \end{itemize}
    \end{solve}
\end{frame}
\resetframecounters%

\begin{frame}{Wieso gibt diese Aufgabe so viele Punkte?}\onslide<+->
    \begin{exercise}
        \begin{enumerate}
            \item[2.] Vereinfachen Sie die folgenden Funktionen so weit wie möglich.
        \end{enumerate}
    \end{exercise}
    \begin{solve}[\(g_1\)]
        \vspace*{-2\abovedisplayskip}\begin{align*}
            g_1 &= (b\wedge (a\vee c))\vee (\overline{b} \wedge (a\vee c))\vee (a\wedge d)
            \only<+->{\\&= ((b\vee \overline{b})\wedge(a\vee c))\vee (a\wedge d)\quad\text{\color{gray}Distributivgesetz}}
            \only<+->{\\&= (1\wedge (a\vee c))\vee (a\wedge d)}
            \only<+->{\\&=(a\vee c)\vee (a\wedge d)}
            \only<+->{\\&= c(a \vee (a\wedge d)) \vee c}
            \only<+->{\\&= a\vee c\quad\text{\color{gray}Absorption}}
        \end{align*}\vspace*{-2\belowdisplayskip}
    \end{solve}
\end{frame}
\resetframecounters%

\begin{frame}{Wieso gibt diese Aufgabe so viele Punkte?}\onslide<+->
    \begin{exercise}
        \begin{enumerate}
            \item[2.] Vereinfachen Sie die folgenden Funktionen so weit wie möglich.
            \item[3.]<3> Wie nennt man die Funktion \(g_2\)?
        \end{enumerate}
    \end{exercise}
    \begin{solve}[\(g_2\)]
        \vspace*{-2\abovedisplayskip}\begin{align*}
            g_2 &= ((a\leftrightarrow b)\vee (a\nleftrightarrow b))\wedge (c\vee \overline{c})
            \only<+->{\\&=1\wedge 1 = 1}\only<+->{\quad\text{\color{gray}Tautologie!}}
        \end{align*}\vspace*{-2\belowdisplayskip}
    \end{solve}
\end{frame}

\end{document}% chktex 17